% chktex-file 44
% chktex-file 8
\documentclass[10pt,leqno]{amsart}
\usepackage{graphicx}
\baselineskip=16pt

\usepackage{indentfirst,csquotes}

\topmargin= .5cm
\textheight= 20cm
\textwidth= 32cc
\baselineskip=16pt

\evensidemargin= .9cm
\oddsidemargin= .9cm

\usepackage{amssymb,amsthm,amsmath}
\usepackage{xcolor,paralist,hyperref,titlesec,fancyhdr,etoolbox}
\newtheorem{theorem}{Theorem}[]
\newtheorem{definition}[theorem]{Definition}
\newtheorem{example}[theorem]{Example}
\newtheorem{lemma}[theorem]{Lemma}
\newtheorem{proposition}[theorem]{Proposition}
\newtheorem{corollary}[theorem]{Corollary}
\newtheorem{conjecture}[theorem]{Conjecture}

\newcommand{\set}[1]{\left\{#1\right\}}
\newcommand{\abs}[1]{\left|#1\right|}
\newcommand{\norm}[1]{\left\|#1\right\|}
\newcommand{\paren}[1]{\left(#1\right)}
\newcommand{\Int}[1]{\mathbb Z / #1 \mathbb Z}
\newcommand{\Unit}[1]{(\Int{#1})^\times}
% \titleformat{\section}[display]{\normalfont\huge\bfseries\centering}{\centering\chaptertitlename\thechapter}{10pt}{\Large}
% \titlespacing*{\section}{0pt}{0ex}{0ex}





\newcommand{\sorry}{\text{sorry}}
\hbadness=99999







\begin{document}

\title{Twelve Lonely Runners}
\author[T. Sungkawichai]{Touch Sungkawichai}
\date{\today}
\address{London, United Kingdom}
\email{touchsungkawichai@gmail.com}
\maketitle
% \let\thefootnote\relax
% \footnotetext{MSC2020: Primary 00A05, Secondary 00A66.}

\begin{abstract}
  The Lonely Runner Conjecture asserts that for any number of runners moving at distinct constant speeds around a unit circle
  from a common starting point, each runner is lonely at some time. A runner is said to be lonely if its distance from every 
  other runner exceeds $1/(k+1)$, where $k+1$ is the total number of runners. A computational method derived by 
  Rosenfeld~\cite{Rosenfeld} and Trakulthongchai~\cite{Trakulthongchai} verified the conjecture for up to ten runners
  In this work, we prove that the Lonely Runner Conjecture in the 12 runners cases holds provided that  
  it holds in the 11 runners case.
\end{abstract}

\bigskip
\section*{Introduction}
The Lonely Runner Conjecture, attributed to Wills~\cite{Wills}, states that for any $n > 2$, if $n$ runners are 
running with constant distinct speed around a unit circle, each runner would be eventually \textit{lonely}.
Formally that means for each runner $i$, 
there is a point in time $t$ such that no runner is within distance $\frac{1}{n}$ of runner $i$.
By using the symmetry of the problem, one may fix a runner to be stationary and measure all other speeds relative to it. 
This yields an equivalent formulation as presented in Conjecture~\ref{lonelyconj}.

\begin{conjecture}
  Let $v_1, \dots, v_k$ be positive real numbers. Then there exists a real number $t$ such that
\[
  \norm{t v_i} \ge \frac{1}{k+1}
\quad \text{for all } i = 1,\dots,n,
\]
  \label{lonelyconj}
\end{conjecture}

The conjecture is trivial for 2 and 3 runners. 
However, as the number of runners increases, the complexity grows exponentially.
While showing the proof for the case of six runners, Bohman, Holzman, and Kleitman~\cite{Bohman} showed that it is 
enough to considered integer speeds, reducing the problem from a continuous to a discrete one. 
Moreover, Tao~\cite{Tao} showed that it is possible to computationally verify the conjecture by only considering integer speeds 
up to $n^{O(n^2)}$, turning an infinite problem into a verifiable one. 
Later on, Malikiosis, Santos, and Schymura~\cite{Malikiosis} reduced the search space drastically by showing that it is enough 
to consider integer speeds up to only $\binom{n+1}{2}^{n-1}$.

Building on the result of Malikiosis~\cite{Malikiosis}, Rosenfeld \cite{Rosenfeld} proved the Lonely Runner Conjecture for 8 runners and 
Trakulthongchai~\cite{Trakulthongchai} improved the computational search to assert for 9 and 10 runners.

In this work, we improved the computational method used by Rosenfeld and Trakulthongchai to conclude 
that if the conjecture holds for $11$ runners, it also holds for 12 runners.

\section*{Rosenfeld's and Trakulthongchai's Computational Method}
By Corollary 3 of \cite{Rosenfeld}, derived from the result of Malikiosis, Santos, Schymura \cite{Malikiosis},
if the $k-1$ Lonely Runner Conjecture holds, the $k$ lonely runners can be proven by verifying that there are no
integral speed set $\set{v_1, \dots, v_k}$ that is a counterexample to the conjecture for 
\[ v_1\dots v_k < B_k \quad \text{where} \quad B_k := \paren{\frac{\binom{k+1}{2}^{k-1}}{k}}^k \]

Rosenfeld introduced the idea of attacking the bound by generating multiple small proofs, 
each of which is a proof that if $v_1, \dots, v_k$ is a counterexample, 
then a particular prime $p$ divides $v_1 \cdots v_k$. 
By accumulating enough proofs for a set of primes $P$, it follows that 
\[ \prod_{p \in P} p \;\Big\vert\; v_1\cdots v_k. \]
which shows $v_1 \cdots v_k \ge \prod_{p \in P} p > B_k$

Following Rosenfeld's approach, Trakulthongchai~\cite{Trakulthongchai} defined proper and improper solution sets. 
The following definition~\ref{defimproper} is a reformulation of Definition 3.1 from \cite{Trakulthongchai}.

\begin{definition}
  Let $p$ be a prime greater than $k+1$ and $\gcd(p, l) = 1$.
  Let $\pi: \Int{pl} \to \Int{p}$ be the natural projection. 
  A speed set $v := \set{v_1, \dots, v_k}$ is such that $v_i \in \pi^{-1}(\Unit{lp})$ for all $i$.
  A speed set $v$ is improper modulo $lp$, write $v \in I(k, l, p)$ if both conditions hold: 
  \begin{enumerate}
    \item for any $t \in \Int{lp}$ there exists $i$ such that $\norm{\frac{tv_i}{lp}} < \frac{1}{k+1}$.
    \item for all $i$, $\gcd(lp, v_1, \dots, v_{i-1}, v_{i+1}, \dots, v_k) = 1$.
  \end{enumerate}
  and $I(k, l, p)$ denotes the set of all improper speed sets for fixed $k, l, p$.
  \label{defimproper}
\end{definition}

Intuitively, the condition states that for every time parameter $t$ modulo $p$, there is a runner who fails 
to achieve the required separation distance $\frac{1}{k+1}$ from the stationary runner. 
Thus the set is a potential obstruction to the Lonely Runner Conjecture when considered modulo $p$.

He showed in Lemma 3.3 of \cite{Trakulthongchai} that if the set of improper speed sets
$I(k, l, p)$ is empty for some value of $l$, then $p \mid v_1\cdots v_k$ in any counterexample.

Trakulthongchai then proposed an efficient method that allows \textit{lifting} a candidate solution set modulo $lp$ 
into a candidate solution set modulo $clp$. Allowing computation to use prime factors of $k+1$ to help speed up the process.
For example, with $k=8$, the idea allows splitting the computation into 3 steps, computing $I(k, 1, p)$ then $I(k, 3, p)$ 
and then $I(k, 9, p)$. 

\section*{Main Idea: Quotient modulo $p$}

The framework introduced in \cite{Trakulthongchai} has a bottleneck in the first phase of computation, the 
calculation of $I(k, 1, p)$. Therefore we aimed to exploit the $\Unit{p}$ symmetry of $I(k, l, p)$ to reduce the amount of 
computation needed for $I(k, 1, p)$. Notice that for $l=1$ and a prime modulus $p$, 
the objects of interest are speed sets whose elements lie in the multiplicative group $(\mathbb{Z}/p\mathbb{Z})^\times$.  

In order to facilitate reasoning about speed set, let $(v_1, \dots, v_k)$ be a canonical representation of speed set 
$\set{v_1, \dots, v_k}$ such that two set are equal if and only if the canonical representations equal.
The representation can be chosen arbitrarily. In this work and the algorithm written in \cite{mycode}, 
we use $(v_1, \dots, v_k)$ such that $v_1 < \cdots < v_k$.

\begin{lemma}
  For any improper speed set $(v_1, \dots, v_k)$ in $I(k, 1, p)$, there exists an improper set
  $(1, u_2, \dots, u_k)$ in $I(k, 1, p)$ and a number $u \in \Unit{p}$ such that 
  $(u, uu_2, \dots, uu_k) = (v_1, \dots, v_k)$.
  \label{lemmaprime}
\end{lemma}

\begin{proof}
  Since $v_1$ is a speed in a speed set modulo $p$, $v_1 \in \Unit{p}$. 
  Let $u = v_1$ and $u_i = v_1^{-1} v_i$ for each $i$.  
  Then, $u u_i = v_1 \cdot v_1^{-1} v_i = v_i$. Hence
  \[
    (u, u u_2, \dots, u u_k) = (v_1, v_2, \dots, v_k).
  \]
  
  It remains to verify that $\{1,u_2,\dots,u_k\}$ is also an improper set. 
  Since $l = 1$, $\gcd(p, v_i) = 1$ for any $i$, so the gcd condition holds trivially. 
  Therefore, it is enough to show the first condition.
  To do that, pick any $t \in \Int{p}$ and set $t' = t v_1^{-1}$.
  Since $(v_1, \dots, v_k)$ is improper, there exists some index $i$ such that
  \[ \frac{1}{k+1} > \norm{\frac{t' v_i}{p}} = \norm{\frac{t v_1^{-1} v_i}{p}} = \norm{\frac{tu_i}{p}} \]
  where in particular, $u_1 = 1$. As both conditions in Definition~\ref{defimproper} hold for the set $\set{1, u_1, \dots, u_k}$, 
  the set is improper.
\end{proof}

The lemma shows an equivalent relation of improper sets in $I(k, 1, p)$ under multiplication by units. 
Consequently, when enumerating possible counterexamples modulo $p$ it suffices to consider
only \textit{normalized} sets of the form $\{1,u_2,\dots,u_k\}$.  This reduces the search space by a factor of $p$. 

When lifting $I(k, l, p)$ to $I(k, cl, p)$ by a multiplier $c$, the \textit{shadow}, as described in \cite{Trakulthongchai},
of $I(k, l, p)$ at level $cl$ is the product of component-wise preimage of the natural projection $\pi: \Int{lp} \to \Int{clp}$.
Formally 
\[ \delta_{cl}I(k, l, p) = \bigcup_{v \in I(k,l,p)} \prod_{i=1}^k \pi^{-1}(v_i) \]
where $\pi: \Int{clp} \to \Int{lp}$ is the natural projection. This motivates the following lemma~\ref{lemmacomp} as an extension of 
lemma~\ref{lemmaprime} for $l > 1$.

\begin{lemma}
  Let $\gcd(l, p) = 1$ and $\pi: \Int{pl} \to \Int{p}$ be the natural projection.
  For any improper set $(v_1, \dots, v_k)$ in $I(k, l, p)$, there is an improper set $(u_1, u_2, \dots, u_k)$ in $I(k, l, p)$
  and a number $u \in (\Int{pl})^\times$ such that $\pi(u_1) = 1$, i.e. $u_1 \equiv 1 \pmod{p}$,
  and \[ (uu_1, uu_2, \dots, uu_k) = (v_1, \dots, v_k) \text{ as sets in } \Int{lp} \]
  \label{lemmacomp}
\end{lemma}

\begin{proof}
  Since $\gcd(p, l) = 1$, the restriction $\pi|_{(\Int{pl})^\times} : (\Int{pl})^\times \to (\Int{p})^\times$ 
  is surjective by the Chinese Remainder Theorem. Since $\set{v_1,\dots,v_k}$ is an improper set, $\pi(v_1) \neq 0$, so 
  $\pi(v_1) \in (\Int{p})^\times$. By surjectivity, there exists $u \in (\Int{pl})^\times$ with $\pi(u) = \pi(v_1)$, 
  and pick $u_i := u^{-1} v_i$ for each $i$. 

  This makes $u u_i = v_i$ for all $i$, so 
  $\{u u_1, \dots, u u_k\} = \{v_1, \dots, v_k\}$. Moreover,
  \[
    \pi(u_1) = \pi(u)^{-1}\pi(v_1) = \pi(v_1)^{-1}\pi(v_1) = 1.
  \]

  It remains to show that $\{u_1, \dots, u_k\}$ is an improper set. For the gcd condition, 
  since $u \in (\Int{pl})^\times$ we have $\gcd(u^{-1}, lp) = 1$, so multiplication by $u^{-1}$ 
  preserves all divisibility relations with respect to $lp$. Hence for each $i$,
  \[
    \gcd(u_1, \dots, \widehat{u}_i, \dots, u_k, lp) = \gcd(v_1, \dots, \widehat{v}_i, \dots, v_k, lp) = 1,
  \]
  where the last equality holds since $\{v_1, \dots, v_k\}$ is an improper set.

  For the second condition, pick any $t \in \Int{pl}$ and set $t' = tu^{-1}$ resembling lemma~\ref{lemmaprime}.
  Since $\{v_1, \dots, v_k\}$ is improper, there exists some index $i$ such that
  \[
    \frac{1}{k+1} > \norm{\frac{t' v_i}{pl}} = \norm{\frac{t u^{-1} v_i}{pl}} 
    = \norm{\frac{t u_i}{pl}}
  \]
  which fulfills condition for $\{u_1, \dots, u_k\}$ to be improper.
\end{proof}

The lemma shows that every improper set in $I(k, l, p)$ is a unit multiple 
of an improper set whose first element projects to $1 \in \Int{p}$. In particular, it suffices to enumerate improper 
sets only among those with $\pi(u_1) = 1$ because 
\[
  I(k, l, p) = \emptyset \iff \{ \{u_1, \dots, u_k\} \in I(k,l,p) : \pi(u_1) = 1 \} = \emptyset.
\]

This reduces the amount of computation needed by roughly a factor of $p$ as it is enough to keep track of 
$I(k, l, p)/\Unit{p}$.

In addition to this normalization, I improved the computational approach of
Trakulthongchai by refactoring the implementation to support parallel
execution and by introducing more memory-efficient data structures.  The new
implementation remains faithful to the original algorithmic ideas but is
significantly faster in practice.

\section*{Introducing intermediate sieves}
In this section, we write $I'(k, l, p) := I(k, l, p)/\Unit{p}$.

This idea was developed from the suggestion introduced by Trakulthongchai in section 7 of \cite{Trakulthongchai}.

The cost of lifting $I'(k, l, p)$ to $I'(k, cl, p)$ is proportional to $|I'(k, l, p)| \cdot c^k$,
since each improper set modulo $lp$ has $c^k$ candidate preimages modulo $lcp$ to check.
In particular, a lift with $c = 2$ costs $|I'(k,l,p)| \cdot 2^k$, whereas a lift with $c = 3$ 
costs $|I'(k,l,p)| \cdot 3^k$, making the $c = 2$ step significantly cheaper for large $k$.

Moreover, I observed that lifting with $c = 2$ can reduce the size of the improper set,
i.e. $|I'(k, 2l, p)| < |I'(k, l, p)|$ in practice. This observation suggests that using a $c=2$ lift can minimize 
$\abs{I'(k, l, p)}$ before passing through the more expensive $c=3$ lift.
Performing an extra cheaper step of $c=2$ lift can shorten the overall computation time by reducing the work needed to be done in 
larger mod.

To put numbers to the observation, I have devised two runs to prove the case of $k=8$. The first one follows Trakulthongchai's method 
described in Algorithm 4.1 of \cite{Trakulthongchai}. The second one involves trying to lift with $c=2$ twice and keep the smallest set.
Specifically, the second method can be described as 
\begin{enumerate}
  \item Generate $S_1 := I'(8, 1, p)$
  \item Compute $S_2 := I'(8, 2, p)$ by lifting $S_1$.
  \item If $\abs{S_2} < \abs{S_1}$, compute $S_4 := I'(8, 4, p)$ by lifting $S_2$.
  \item Take $T_1$ as the smallest set among $S_1, S_2$ and $S_4$.
  \item Compute $T_3$ by lifting $T_1$ with $c=3$. 
  \item Compute $T_9$ by lifting $T_3$ with $c=3$.
\end{enumerate}

In the second, the $c=2$ lifts are unnessary as both programs are capable of finding the proof for each prime. However, 
by doing the extra steps, the amount of work left at $c=3$ lift was reduced, shortening the compute time. 
By running the two methods on a 10-core Apple M4 processor (Macbook Air 2025), 
the first method takes 18 seconds while the second methods uses 12 seconds.

This speedup can be dramatic in certain cases, for example in the case of $k=9$. One idea is to generate $S_1 := I'(9, 1, p)$, 
then lift it to $S_2 := I'(9, 2, p)$ with $c = 2$ and lift it to $S_{10} := I'(9, 10, p)$ with $c = 5$. My proposal is 
\begin{enumerate}
  \item Generate $S_1 := I'(9, 1, p)$
  \item Compute $S_2 := I'(9, 2, p)$ by lifting $S_1$ with $c=2$.
  \item Compute $S_4 := I'(9, 4, p)$ by lifting $S_2$ with $c=2$.
  \item Take $T_1$ to be the smaller set among $S_2$ and $S_4$.
  \item Compute $T_3$ by lifting $T_1$ with $c=3$. At this point, 
    $T_3$ can be either $I'(9, 6, p)$ or $I'(9, 12, p)$.
  \item Take $U_1$ to be the smaller set among $T_1$ and $T_3$.
  \item Compute $U_5$ by lifting $U_1$ with $c = 5$.
    At this point, $U_5 = I'(9, l, p)$ for some $l$ such that $10 \mid l$.
\end{enumerate}

For a specific number, the first method without the intermediate step could not compute $I'(9, 10, 53)$ within 15 minutes while it takes 
under 5 seconds with the second method. This opens up the possibility of lifting an improper set to a larger prime number and could help with 
the verification of the conjecture for 11 runners.

\section*{Result}

Rosenfeld wrote \cite{rosenfeldcode} to verify the conjecture for eight runners and Trakulthongchai developed \cite{tanupatcode} to 
verify the conjecture of nine and ten runners. His computation takes 
15 minutes for $k=8$ and approximately 23 hours for $k=9$. Building on the framework of \cite{tanupatcode}, I developed \cite{mycode},
with assistance from AI-based tools, to improve performance and scalability.

Using the improvements proposed, I managed to complete the proof of $k=8$ in 12 seconds, $k=9$ in 181 seconds, and $k=11$ in approximately 
\sorry hours on a 10-core Apple M4 processor (Macbook Air 2025). The source of the computation can be found in \cite{mycode}

For the twelve runners case, let \[ P_{11} := \set{ \begin{aligned} 
      23, 131, 137, 139, 149, 151, 157, 163, 167, 173, 179, 181, 191, \\ 
      193, 197, 199, 211, 223, 227, 229, 233, 239, 241, 251, 257, 263, \\
      269, 271, 277, 281, 283, 293, 307, 311, 313, 317, 331, 337, 347, \\
      349, 353, 359, 367, 373, 379, 383, 389, 397, 401, 409, 419, 421, \\
      431, 433, 439, 443, 449, 457, 461, 463, 467, 479, 487, 491, 499, \\ 
      503, 509, 521, 523, 541, 547, 557, 563, 569, 571, 577 \\
    \end{aligned}
  } \]

Using the algorithm described in the previous section, I have verified
that for every prime $p \in P_{11}$ there exists $l \in \set{9, 18, 36}$ such that the set $I(k, l, p)$ is empty.
Equivalently, any counterexample to the $12$-runner case of the
Lonely Runner Conjecture must have speeds divisible by each prime in
$P_{11}$. Consequently, if such a counterexample exists, its product of speeds must be
divisible by the product $\prod_{p \in P_{11}} p$

\begin{theorem}
  The Lonely Runner Conjecture for $k=11$ holds if the Lonely Runner Conjecture for $k=10$ holds. 
\end{theorem}
\begin{proof}
  If the conjecture holds for eleven runners, then the result of \cite{Rosenfeld} asserts that it suffices 
  to check integral solutions up to $B_{11}$.

  As $B_{11}$, the bound needed to reduce the problem of twelve runners to eleven runners, is less than $e^{434.5}$ and 
  \[ \prod_{p \in P_{11}} p > e^{435} > e^{434.5} > B_{11} = \paren{\frac{\binom{k+1}{2}^{k-1}}{k}}^k \] 
  There is no counterexample to the 12 runners case.
\end{proof}

\section*{Further work}
The case of 11 runners remains unsolved. Since 11 is prime, the bottleneck of the computation is not at finding $I(k, 1, p)$ but at 
lifting $I(k, 1, p)$ to $I(k, 11, p)$. Similar arguments may potentially speed up calculation but will require more careful consideration.

\section*{Acknowledgement}
I am grateful for Tanupat Trakulthongchai for his help and guidance throughout the preparation of this article.

\bibliographystyle{plain}
\bibliography{ref}

\end{document}
